% A neutral six-page deck, written for the thesis figure of the reading view.
% Its text is our own; it contains no course material and no names. Upload the
% built PDF to a Library, open it in the reading view, and take the
% screenshots there; the figure's caption says which views they show.
% Build: latexmk -pdf -outdir=../../build sample-deck.tex
\documentclass[aspectratio=169,11pt]{beamer}
\usetheme{default}
\setbeamertemplate{navigation symbols}{}
\setbeamertemplate{footline}[frame number]
\usepackage[T1]{fontenc}
\usepackage{lmodern}
\usepackage{booktabs}

\title{Binary Search}
\subtitle{Finding an item in a sorted list}
\date{}

\begin{document}

\begin{frame}
  \titlepage
\end{frame}

\begin{frame}{The idea}
  \begin{itemize}
    \item Look at the middle element of the part that is left.
    \item If it is the item, stop.
    \item If the item is smaller, continue in the left half;
          otherwise continue in the right half.
    \item Every step halves the part that is left.
  \end{itemize}
\end{frame}

\begin{frame}{Step by step}
  Find 23 in the sorted list 3, 8, 12, 17, 23, 31, 42.
  \medskip
  \begin{center}
  \begin{tabular}{@{}cccl@{}}
    \toprule
    low & high & middle & comparison \\
    \midrule
    0 & 6 & 3 (17) & 23 > 17: go right \\
    4 & 6 & 5 (31) & 23 < 31: go left \\
    4 & 4 & 4 (23) & found \\
    \bottomrule
  \end{tabular}
  \end{center}
  \medskip
  Three comparisons instead of five.
\end{frame}

\begin{frame}{Cost}
  \begin{itemize}
    \item A list of $n$ items needs at most $\lfloor \log_2 n \rfloor + 1$ comparisons.
    \item 1,000 items: at most 10 comparisons.
    \item 1,000,000 items: at most 20 comparisons.
    \item A linear search may need all $n$.
  \end{itemize}
\end{frame}

\begin{frame}{Preconditions and pitfalls}
  \begin{itemize}
    \item The list must be sorted, or the halves mean nothing.
    \item With fixed-width integers, $(\mathit{low} + \mathit{high})/2$ can overflow;
          $\mathit{low} + (\mathit{high} - \mathit{low})/2$ cannot.
    \item Loop bounds are a common source of off-by-one errors.
  \end{itemize}
\end{frame}

\begin{frame}{Summary}
  \begin{itemize}
    \item Halve the part that is left at every step.
    \item At most about $\log_2 n$ comparisons.
    \item Works only on sorted input.
  \end{itemize}
\end{frame}

\end{document}
